\begin{problem}[34.10]
  Factor the polynomial $4x^2 - 4x + 8$ into a product of irreducibles viewing it as an element of the integral domain $\Z[x]$; of the integral domain $\Q[x]$; of the integral domain $\Z_{11}[x]$.
\end{problem}

\begin{solution}
\end{solution}

\begin{problem}[34.25]
  Prove that if $p$ is a prime in an integral domain $D$, then $p$ is an irreducible.
\end{problem}

\begin{solution}
\end{solution}

\begin{problem}[34.26]
  Prove that if $p$ is an irreducible in a UFD, then $p$ is a prime.
\end{problem}

\begin{solution}
\end{solution}

\begin{problem}[34.27]
  For a commutative ring $R$ with unity show that the relation $a \sim b$ if $a$ is an associate of $b$ (that is, if $a = bu$ for $u$ a unit in $R$) is an equivalence relation on $R$.
\end{problem}

\begin{solution}
\end{solution}

\begin{problem}[34.29]
  Let $D$ be a UFD. Show that a nonconstant divisor of a primitive polynomial in $D[x]$ is again a primitive polynomial.
\end{problem}

\begin{solution}
\end{solution}

\begin{problem}[34.30]
  Show that in a PID, every proper ideal is contained in a maximal ideal. [Hint: Use Lemma 34.10.]
\end{problem}

\begin{solution}
\end{solution}

\begin{problem}[35.2]
  The function $\nu$ for $\Z[x]$ given by $\nu(f(x)) = (\text{degree of}~f(x))$ for $f(x) \in \Z[x]$, $f(x) \ne 0$.
\end{problem}

\begin{solution}
\end{solution}

\begin{problem}[35.8]
  Following the idea of Exercise 6 and referring to Exercise 7, express the positive gcd of 49,349 and 15,555 in $\Z$ in the form $\lambda(49,349) + \mu(15,555)$ for $\lambda, \mu \in \Z$.
\end{problem}

\begin{solution}
\end{solution}

\begin{problem}[35.9]
  Find the gcd of
  \[%
    x^{10} - 3x^9 + 3x^8 - 11x^7 + 11x^6 - 11x^5 + 19x^4 - 13x^3 + 8x^2 - 9x + 3
  ,\]%
  and
  \[%
    x^6 - 3x^5 + 3x^4 - 9x^3 + 5x^2 - 5x + 2
  ,\]%
  in $\Q[x]$.
\end{problem}

\begin{solution}
\end{solution}

\begin{problem}[35.15]
  Let $D$ be a Euclidean domain and let $\nu$ be a Euclidean norm on $D$. Show that if $a$ and $b$ are associates in $D$, then $\nu(a) = \nu(b)$.
\end{problem}

\begin{solution}
\end{solution}
